% ============================================================
%  09_conclusions.tex  –  Chapter 7: Conclusions and Future Work
% ============================================================

\chapter{Conclusions and Future Work}
\label{ch:conclusions}

\section{Conclusions}
\label{sec:conclusions}

In conclusion, this thesis presents the indoor air station (\airnest{}), a complete, functional
low-cost IoT-based environmental monitoring system designed for residential use. For the part
related to taking readings it features MQ gas sensors, a DHT11 sensor and a KY-037 sensor. The
Raspberry Pi 4B is used as the central processing unit. The project embodies a system that
collects, stores, computes and presents multiple air quality parameters in real life.

The hardware product reads: five gas sensors: carbon monoxide, hydrogen, LPG/butane,
alcohol/benzene; temperature, humidity; ambient and sudden spikes when it comes to the sound
level (at 30 minutes intervals). The raw ADC values are converted to voltages, corrected for
the voltage divider configuration. Afterwards the power-law formula to compute PPM
concentrations is used. In order to have an increased accuracy, a temperature and humidity
correction factor derived from Hanwei application notes~\cite{hanwei} is used for the PPM
values. Warnings are generated when the values go over the approximate safety thresholds.

The data related to the readings is locally saved in a CSV file (on the Raspberry Pi). This
way a backup is always present, in case something stops working properly. Each reading is also
sent, as a JSON payload over HTTP, to the Flask server, which is hosted on AWS EC2. There, the
data is inserted in the PostgreSQL database (hosted on AWS RDS). The cloud based architecture
makes sure that the data is always available, even if the station is stopped.

The website provides a complete interface for interacting with the data. All the collected and
computed readings can be seen on the Data page, through the line charts (rendered using
Chart.js~\cite{chartjs}). The user authentication system supports registration with email
verification. The password is stored in a secure manner. For this bcrypt hashing and a
token-based password reset flow were used. Users who are authenticated can download all the
sensor related data as a CSV file. They can also choose to download the mobile app (through
\texttt{.apk}), which offers the users the possibility to look at the data charts and
predictions, as well as access to write in the contact form. The app is made from the already
existing website pages (home page, data page, prediction page and form page). The website also
has an admin page (accessed through the admin account) from which the logged in users can be
seen and also the contact messages. This part also has a toggle to put the entire website in
``maintenance'' mode. The contact form allows users to send messages to the administrator,
even if they are not registered or logged in.

This product achieves its intended purpose: to provide a low-cost, open and customizable
apartment air quality monitoring station, whose data can be always accessed online. It also
gives the user full data visualization and ownership capabilities.

\section{Future Work}
\label{sec:future-work}

Several directions for further development have been identified during the course of this
project, ranging from immediate hardware and software improvements to planned new features.

The DHT11 sensor should be moved even further away from the rest of the cramped hardware parts,
since the heat radiating from the rest of the components could make it register wrong readings
instead. The best solution would be to move it to a new breadboard, at least 20 centimeters
away from the rest.

The MQ7 sensor was found to read values higher than the expected amount. This is due to the
fact that its accuracy tolerance is of around $\pm$15--30\,\% even when properly calibrated and
this sensor also being more accurate for higher value readings. In the future it would be best
to replace this sensor with a more precise one when it comes to lower values.

The neighbourhood label is currently configured through a terminal prompt at startup and cannot
be changed while the station is running. A more practical solution would be to read this value
from a configuration file at startup, allowing the user to update it between sessions without
modifying the code.

For further training the AI models, dataset size must be bigger in order to provide better
results. The local sensor history is still limited. Model accuracy will improve as long as more
data is collected and used for retraining.

Adding sensors specific for reading values for PM2.5 and PM10 would further improve the
capabilities of the station. This way it would be brought closer to the feature set of
commercial air quality monitors reviewed in Chapter 2.

Placing the hardware station in a protective case would increase the likelihood of the project
to be compared more to a real, commercial product, besides offering protection to the
components.

A cooling mechanism should be added for the entire hardware part. The Raspberry Pi overheats
after a time of continuous operation and the MQ sensors need a preheating time to give accurate
reading, time which takes longer the warmer it is around.

Adding the option to start the Raspberry Pi script through the mobile application by using
Bluetooth would bring it closer to a real, sellable product, while also making accessibility
easier for the user.
